\chapter{A Biofeedback Framework for Attention and Memory}

Functional neuroimaging techniques have transformed neuroscience studies by providing researchers the opportunity to catch a glimpse of the brain from the outside (\cite{poldrack2018new}, p. 2). Despite the range of spatially and temporally resolved signatures of neural activity that various neuroimaging techniques supply, these signals typically constrain conclusions drawn from brain and behavior associations to be correlational rather than causal evidence for a directional relationship (\cite{atkinson2005visual}, cf. \cite{weber2010functional}). Tight correlations between neural signals and behavior are nevertheless important for guiding future research towards causal investigations of the brain on behavior. Studies aiming to elucidate further understanding of the underlying neural mechanisms subserving human cognitive control function and behavior must carefully consider designs where causal interventions on brain signals are minimally invasive and not harmful to human subjects. One emerging approach is the use of real-time closed-loop behavior-interventions with engineered biofeedback (\cite{laconte2011decoding, alkoby2018can}).

Previous engineered biofeedback implementations have used real-time continuous neurofeedback (NF) signals to reward particular signatures of brain activity, thus `training' participants to modify their brain activity on subsequent trials (\cite{alkoby2018can}). NF has been utilized for behavioral interventions in healthy adults (\cite{hoedlmoser2008instrumental, zoefel2011neurofeedback, wang2013neurofeedback, gruzelier2014eegA, gruzelier2014eegB}) as well as in treatments for patients with a gamut of disorders, including stroke (\cite{rayegani2014effect, doppelmayr2007attempt}), epilepsy (\cite{kotchoubey2001modification, walker2005neurofeedback, sterman2006foundation, tan2009meta, strehl2014sustained}), and attention deficit hyperactivity disorder (ADHD; \cite{thompson2005neurofeedback, leins2007neurofeedback, arns2009efficacy, sonuga2013nonpharmacological}). Though, most closed-loop NF approaches have been geared at brain training interventions, where the efficacy of these training treatments is apparent in some but absent in others (\cite{alkoby2018can, weber2011predicting}). Yet, understanding of NF training inefficacy --- and the characteristics driving it --- are limited as many studies generally fail to report on individual variation within their samples (\cite{alkoby2018can}).

Considering these limitations, I aim to elucidate the efficacy of real-time biofeedback approaches for causally manipulating an experimental paradigm. Specifically, I will discuss the role that pupillometry can serve as a temporally resolved toolkit within a feasible and non-invasive conceptual biofeedback implementation for closed-loop interventions to causally link lapses in sustained attention with poorer memory performance. Here, four conceptual methodological implementations will be described, including a real-time attention-orienting triggering procedure using tonic pupil diameter on a trial-by-trial basis, as well as a predictive modeling approach using machine learning trained on data from each individual's performance on an initial episodic memory task; this participant-specific model can then be implemented to causally reorient attention --- using feedback cues manipulated as a function of predicted lapses of sustained attention --- with the ultimate goal of limiting the amount of trials where memory performance suffers during retrieval.

\section{Sustained Attention and Goal-Directed Memory}
Sustained attention is fundamental for successful goal-directed episodic memory (\cite{baddeley1984attention, craik1996effects, anderson1995status, evans2019pre, curran2004effects}). Lapses in attention are correlated with higher likelihoods of memory failures in goal-directed behavioral tasks (\cite{madore2020memory}). Both task-based (\cite{esterman2013zone, rosenberg2013sustaining}) and self-report questionnaire (\cite{ophir2009cognitive, ralph2014media, kessler2005world, patton1995factor, green2007action, carriere2008everyday, carriere2013wandering}) measures of sustained attention have informed critical relations between trait-level attention lapses and memory performance (\cite{madore2020memory}). Laboratory-based assays of spontaneous attention lapses often leverage pre-goal tonic posterior alpha power and pupillometry signals, given prior work demonstrating relations between trait-level inattention and both tonic alpha power and variability in tonic pupil diameter (\cite{unsworth2017importance, klimesch2012alpha}). Electroencephalography (EEG) and pupillometry measurements of attention (\cite{macdonald2011trial, unsworth2018tracking}) further inform negative relations between trait-level posterior alpha power and pupil diameter, respectively, with remembering and forgetting (\cite{madore2020memory}). Despite the useful insights drawn from correlations between neuroimaging readouts and behavioral measures, additional explanatory progress in elucidating the relationship between attention and memory (i.e., both working and episodic forms) could come from the utilization of neuroimaging toolkits with engineered biofeedback in experimental paradigms. In this context, engineered biofeedback can be used to directly manipulate attention during experimental memory tasks in real-time. Leveraging real-time readout from neuroimaging techniques, which are typically used to generate correlations between neural signatures and behavior, provides an underappreciated opportunity to draw causal inferences from methods that are usually described as being correlational in nature.

Behavioral assays of sustained attention, independent of the primary experimental task, also reveal that trait-level variation in sustained attention maps onto individual differences in forgetting (\cite{madore2020memory}). One such measure is the  gradual-onset continuous performance task (gradCPT; \cite{esterman2013zone, rosenberg2013sustaining}), which is comprised of two metrics --- commission error rate and reaction time (RT) variability. These assays of attention, although informative, limit the extensibility of the conclusions that can be drawn from their correlational relations with memory failures. These methodological and analytical limitations supply fruitful opportunities for memory researchers to explore the utilization of closed-loop designs where fluctuations of real-time signals --- exceeding some threshold on a trial-by-trial basis during a task, for example --- directly influence which level of the independent variable is utilized on the subsequent trial (\cite{laconte2011decoding}). Such a dynamic, biofeedback-inspired, non-invasive approach will give researchers the ability to draw more concrete causal inferences for both basic science aims and translational interventions.

\section{Real-Time Closed-Loop Biofeedback Interventions}
Closed-loop implementations sufficient for testing real-time causal attention lapse effects on memory are limited. Broadly, closed-loop approaches are generally constrained either spatially or temporally --- and may be more feasible and/or less invasive --- depending on the toolkits and techniques used. With a completely non-invasive approach, deBettencourt et al. (\cite*{debettencourt2019real}) implemented a real-time triggering procedure investigating causal effects of lapses in attention on working memory by assessing fluctuations in RT across trials. Although these findings support a direct influence of attention on upcoming working memory, this approach uses a coarse behavioral index of attention state that might be better captured by neural signatures strongly associated with sustained attention fluctuations (\cite{murray2011markers, myers2014oscillatory, adam2015contribution, giesbrecht2006pre, debettencourt2019real}).

\subsection{Functional Magnetic Resonance Imaging}
deBettencourt et al. (\cite*{megan2015closed}) have explored the use of real-time functional magnetic resonance imaging (rtfMRI; \cite{weiskopf2004principles, laconte2011decoding, sulzer2013real}) with multivariate pattern analysis (MVPA; \cite{norman2006beyond}). With a continuous feedback signal reflecting moment-to-moment fluctuations in sustained attention, they used a closed-loop intervention where participants were provided with a customized feedback signal of the difference between the amount of task-relevant to task-irrelevant information active in their brains on a go/no-go task for image categories. Both feedback and control participants knew about the attention-feedback manipulation that was in place (i.e., that the task would be easier if they were paying attention and harder if they were not). The critical difference between the groups was that the feedback group reacted to \emph{and} drove subsequent feedback from their own blood-oxygen-level-dependent (BOLD) signal, while the control group received sham feedback yoked from other participants’ brain signals. The other participants’ brain signals that were used were matched in age, gender, and handedness to each of the control participants. After the experiment concluded, all participants were asked the following: \say{Did you feel that you could control the image with your brain?} Their results showed that the sham feedback impacted the control participants’ attention, despite only 4 out of 16 control participants self-reporting that they had some feeling of brain control (11 out of 16 participants reported such in the experimental feedback group). Importantly, behavioral performance for control participants whose feedback was yoked to the matched experimental group participants did not improve from pre- to post-training, unlike the results observed in the experimental feedback group; this approach improved participants' behavioral performance after one training session. 

A closed-loop system with rtfMRI thus provides a spatially resolved feedback signal. Spatially resolved signals are especially useful for measuring cognitive processes that are distributed broadly throughout various brain regions and networks. To further illustrate, deBettencourt et al. (\cite*{megan2015closed}) used the MVPA approach in real-time to decode attention using a whole-brain BOLD signal to ascertain a participant’s attention to the task-relevant versus the task-irrelevant stimuli. While this spatially resolved technique is useful for distributed cognitive processes like attention, a closed-loop sustained attention reorienting manipulation using real-time pupillometry feedback would be less spatially resolved, but more temporally resolved \emph{and} feasible to implement. Using real-time tonic pupil diameter signals to trigger attention reorienting interventions for memory has the potential to scale tremendously. An rtfMRI approach generally relies on expensive scanning equipment, costly hourly rates to run, and practical limitations with respect to who can get into the scanner for a variety of health and safety reasons, thus drastically limiting the extensibility of investigating the causal links between lapses in sustained attention on memory performance across a wide gamut of time scales (i.e., from minutes to hours), across a majority of the entire lifespan, and even in populations who would be put at substantial risk in the MR environment. 

Decoding neural signatures in real-time to assess attentional states with rtfMRI provides a more spatially resolved signal where the key brain regions and neural mechanisms involved in distinct representations of attention can be more clearly elucidated. Furthermore, the rtfMRI approach supplies more detailed insight into the specific information a participant is attending to (\cite{megan2015closed}). While this approach has not been used to specifically investigate the downstream influence of attention lapses on working memory, it would be highly informative to the field; however, the feasibility and invasiveness of the closed-loop rtfMRI approach makes it challenging to implement at large scales and for individuals who may be challenging to scan (e.g., young children, individuals with claustrophobia, older adults with health conditions).

\subsection{Pupillometry and RT}
Pupillometry --- or the continuous measure of pupil diameter --- is a technique that provides a window into the brain through the eyes and has transformed understanding of mental states (i.e., shifts in attention and perception) in psychological and cognitive science research (\cite{laeng2012pupillometry, hess1960pupil, hess1964pupil, kahneman1966pupil}). Neuroscience research has revealed a strong relationship between the noradrenergic (NE) system's locus coeruleus (LC) and pupil dilation, and the LC-NE neuromodulatory system tightly regulates task engagement by way of increased neuronal gain in task-relevant regions (\cite{laeng2012pupillometry, murphy2011pupillometry, wilhelm1999pupillography, aston2005integrative, berridge2003locus, aston1981nonrepinephrine, sara2009locus, foote1980impulse}). The fine temporal resolution of the continuous pupillometry signal can pair well with spatially resolved neuroimaging approaches like fMRI, which, to illustrate, has recently uncovered new understanding of the neural mechanisms underlying subprocesses of working memory load and stimulus salience (\cite{fietz2021pupillometry}). Thus, using a temporally resolved, continuous biofeedback signal (i.e., pupil diameter) that is tightly correlated with the LC-NE neuromodulatory system --- and subsequently attention state --- is a feasible approach for minimally invasive closed-loop interventions on reorienting in-the-moment attention lapsing to boost subsequent memory.

Tonic pupil diameter could therefore be a potentially informative biofeedback signal for intervening moments where attention lapses are most likely to occur and impact subsequent learning and remembering. Tonic pupil diameter is the sustained (or basal) component of pupillary response, which contrasts the phasic pupil response, or the transient component (\cite{beatty1982phasic, peysakhovich2017impact}). Tonic pupil diameter is measured as a function of time for each trial, and this basal component of pupillary response occurs prior to each phasic period of pupillary response to the task-relevant stimuli. Over the course of each trial, this tonic pupil diameter signal will be used to determine lapses in sustained attention. Previously, Madore et al. (\cite*{madore2020memory}) assayed spontaneous attention lapses with trial-to-trial retrieval data using tonic pupil diameter from the 1 s directly preceding a retrieval goal cue and a subsequent object probe in an episodic memory task. To obtain trial-by-trial tonic pupillary response at retrieval for pre-goal and pre-probe cues, they used an 8-s intertrial interval (ITI) before the appearance of the 2-s goal cue, which was then followed by an 8-s interstimulus interval before the appearance of the 2-s object probe. Extracting pupil diameter in the 1 s directly before the goal and retrieval cue appearances provides a sufficient signature of trial-level spontaneous tonic lapses (\cite{madore2020memory}). This experimental design with trial-by-trial pupillary readout is critical for computing attentional state across a multi-trial period. 

A similar multi-trial approach has been used with behavioral RT in prior work by deBettencourt and colleagues (\cite*{debettencourt2019real}), where mean RT was assessed over \emph{n - 3} preceding trials and compared to a bidirectional threshold that was the dynamically updating sum/difference of the cumulative mean and standard deviation for all preceding trials (\emph{1, 2, ..., i}). Though, unlike the pupil measures which occur directly before the behavior, RT is a part of the memory retrieval behavior and is thus measured directly following the moments leading up to memory retrieval where tonic pupil diameter is measured. In this coarse behavioral approach (i.e., \cite{debettencourt2019real}), mean recent RTs which drift outside the interval defined by the dynamically updating trailing mean RT plus-or-minus the standard deviation of trailing RT --- for each individual --- would signal attention lapses. This is because faster RTs than usual on a continuous performance task (\cite{megan2015closed, debettencourt2018forgetting, robertson1997oops, rosenberg2013sustaining}) are suggestive of attention lapsing by way of higher error rates on infrequent trials (\cite{megan2015closed, debettencourt2018forgetting, esterman2013zone}). Taking this approach into consideration, the trailing mean of the 1-s pre-cue tonic pupil diameter measures described in Madore et al. (\cite*{madore2020memory}) can thus be used to assay fluctuations in sustained attention over the immediately preceding 3 trials for a closed-loop biofeedback intervention implementation with real-time pupillometry.

Whether a closed-loop intervention uses pupil size or RT variability across \emph{n} preceding trials, the real-time signal of attention only measures the \emph{level of attention}, and not what is actually being attended to --- unlike the spatially resolved rtfMRI approach with voxel pattern decoding described previously (\cite{megan2015closed, debettencourt2019real, rosenberg2013sustaining, yoss1970pupil}). Unlike rtfMRI, the temporally resolved pupillometry signal is particularly useful if the overarching goal of the closed-loop intervention is concerned with detecting moments of attention lapses \emph{independent} of the context of that attention lapse. In the conceptual framework proposed here, I am particularly concerned with identifying any moment where attention is lapsing in preparatory goal-directed periods prior to episodic memory retrieval. Because the anticipated behavioral memory paradigms would be designed to focus on a single encoding and retrieval memory task, decoding what participants are specifically attending (or not attending) to is not of particular interest here, and thus would not require a rtfMRI implementation.

Closed-loop implementations designed around pupil diameter or RT variability are naturally constrained by the scale at which the intervention occurs at. Specifically, variability is defined over multiple measurements, typically at the trial-level. Thus, pooling these signals over multiple trials to assess variability is one way to make a quantitative decision about attention lapsing. For instance, past work that has used significant deviations of RT --- beyond specified thresholds --- to signal lapses in sustained attention were able to do so by gathering RT signal over three trials prior to the current trial (\cite{debettencourt2019real}). Although this real-time triggering procedure is informative for understanding causal links between attention and memory, an approach using signal deviations at a scale of \emph{n - 3} trials is a coarse behavioral assay limiting how confidently such an intervention can directly speak to temporally-resolved moment-to-moment lapses of attention (see \cite{madore2020memory}). Thus, in addition to using slow rhythm (i.e., this 3-trials-back approach) measures of sustained attention, I aim to design and discuss a conceptual framework which includes a finer grained assessment of a participant's lapsing attention state. This approach will also take into account individual variability, therefore serving as a more careful measurement of behavioral signal with respect to individual differences that could potentially reduce the signal-to-noise ratio (see \cite{alkoby2018can}). Below, I outline four potential methodological concepts that could potentially resolve these considerations, albeit each with their own unique sets of inherent strengths and limitations.

\section{Proposed Closed-Loop Intervention on Attention to Causally Manipulate Memory Performance}
Overall, a closed-loop real-time intervention approach for reorienting in-the-moment attentional lapsing is a sensitive measure for understanding the causal role of attention on memory in the preparatory periods just prior to learning and remembering (\cite{megan2015closed}). This contrasts a simpler, more traditional experimental design with, for example, an \say{attention lapsing} and a \say{control} condition where attention lapses are induced by some behavioral mechanism; this coarse behavioral approach would not account for in-the-moment variability of individuals' attentional abilities (\cite{drew2008neural}) nor would it control for instances where someone is actually paying attention on a supposed \say{attention lapsing} trial. Using a continuous biofeedback signal built upon each individual's own biofeedback and attentional readout is therefore essential for the real-time capturing \emph{and} reorienting of attention lapses in the tonic preparatory periods crucial for learning and memory (\cite{madore2020memory}). 

Here, I lay out the initial conceptual considerations and practical requirements necessary for the successful development of closed-loop intervention methodology using pupillometry techniques. I particularly focus this aim with the goal of causally manipulating goal-directed episodic memory on behavioral paradigms with distinct encoding and retrieval phases. This framework could apply to assessing other forms of memory, such as working memory through a whole report precision task (\cite{adam2017clear}), for instance. The design principles with respect to causally manipulating preparatory moments of lapsing attention could be applied to other NF closed-loop systems, such as with EEG and fMRI; however, this is outside the scope of the current discussion. The practical benefits of pupillometry, alongside the limitations of other neuroimaging closed-loop approaches, motivate the conceptual framework discussed throughout the rest of this paper.

\subsection{Concept 1: Threshold-Based Tonic Pupil Diameter Variability Method over \emph{n - 3} Preceding Trials}
\label{section:concept-1}
This first approach would require a task design that could, in real-time, compute variability in tonic pupil diameter over the immediately preceding three trials. When tonic pupil variability is significantly greater than the experiment-wide trailing sum of the mean and standard deviation of tonic pupil diameter, this would trigger an attention reorienting cue (see Figure \ref{fig:key-figure-trial-by-trial}a). This is similar to work by deBettencourt et al. (\cite*{debettencourt2019real}) described previously, where they assessed RT variability for real-time triggering of attention cues; mean RT was assessed over \emph{n - 3} preceding trials and compared to a bidirectional threshold that was the dynamically updating sum/difference of the cumulative mean and standard deviation for all preceding trials (\emph{1, 2, ..., i}). In this particular setup, if the trailing mean RT for the \emph{n - 3} preceding trials was greater than the cumulative sum threshold, then a slow memory probe was triggered to dampen faster forms of behavior, while if the trailing mean RT was less than the cumulative difference threshold, then a fast memory probe was triggered to increase slower forms of behavior. On the contrary, if the training mean RT was within the threshold --- or if the performance was inaccurate on the preceding trials --- then the task jumped directly to trial \emph{i + 1}.

Regardless of the method used to compute the threshold for triggering the attention intervention, if mean tonic pupil variability for the preceding \emph{n - 3} trials exceeds this threshold, then the experimental paradigm should automatically trigger a cue that would aim to reorient attention before continuing the task. In the context of episodic memory retrieval, when pupil diameter exceeds this threshold, the task could either implicitly direct attention back towards the primary retrieval task or even provide explicit feedback to participants indicating that their attention is wandering. The specifics regarding the type of attention reorienting cue used depends on the particular memory paradigm and research question. The important consideration here relies on whether or not there should be a bidirectional threshold (i.e., an upper \emph{and} lower bound) like the RT variability metric described above (\cite{debettencourt2019real}). In principle, creating an interval band to constrain tonic pupil diameter would have transformative downstream potential for training individuals to modulate their attention state via their pupil diameter --- especially by being able to increase pupillary constriction \emph{or} dilation to stay within their particular range for maintaining optimal sustained attention for memory performance. The strength of using a dynamically updating threshold band would be that this interval takes into account individual differences, as opposed to a fixed, pre-determined threshold approach that may not adequately reflect the appropriate pupil diameter cutoff for all participants. Either way, a weakness of this approach concerns whether or not a dynamically updating cumulative sum/difference of the mean and standard deviation of tonic pupil diameter is sufficient to accurately identify momentary lapses of sustained attention.

\begin{figure}[ht!] 
    \centering
    \includegraphics[width=1\linewidth]{figs/n-3_trial_tonic.png}
    \includegraphics[width=1\linewidth]{figs/model_trial-by-trial-label.png}
    \caption[Short title for List of Figures]{\textbf{\textit{a,}} Threshold-based detection of attention lapses with pupillometry over immediately preceding \emph{n - 3} trials. \textbf{\textit{b,}} Machine learning predictive trial-by-trial attention lapses with pupillometry and reaction time. Created with \href{BioRender.com}{BioRender.com}.}
    \label{fig:key-figure-trial-by-trial}
\end{figure}

\subsection{Concept 2: Trial-Level (\emph{n - 1}) Predictive Model-Based Pupil Diameter × RT Method}
\label{section:concept-2}
This second approach varies from \emph{Concept 1} (Section \ref{section:concept-1}) in that it would not rely directly on dynamically updated or fixed thresholds to trigger causal interventions on attention and subsequent memory. Instead, participants would engage in an initial training period where they would complete the behavioral memory task without any biofeedback-based intervention (see Figure \ref{fig:key-figure-trial-by-trial}b). This initial period would be used to gather participant-specific training data, which would subsequently be used to train a predictive model with machine learning techniques. This individual-specific model would predict the probability that a particular trial would be correct (1) or incorrect (0) based on a parameterized fit taking into account mean pupil diameter and RT for the trial immediately preceding it (\emph{n - 1}). To account for potential model fitting convergence issues or poor predictive power, a pre-determined set of possible models could all be fit and compared to select the model that best predicts each participants' individual-level performance. 

Once the appropriate model is validated, the participant would then complete the task again with cues to reorient attention. Attention lapses in this case would be determined via tonic pupil diameter and RT for the \emph{n - 1} trial immediately preceding the current trial. These values would be tested against the model, and a value of 0 or 1 would be used to indicate whether or not the participant would be likely to fail on the current trial. Attention reorienting cues would then intervene depending on the probability that the individual's particular tonic pupil and RT for the previous \emph{n - 1} trial suggested that performance is likely to suffer on the current trial. Further, to quantify the effectiveness of the attention reorienting intervention, a within-subjects control manipulation will be implemented. Specifically, on half of the trials where an attention lapse is detected, participants will receive the reorienting intervention, while on the other half they will not. To avoid potential global cueing effects on performance, the attention reorienting intervention will be included on a random subset (e.g., 25\%) of trials where participants are attending. This value will depend on the ratio of attention lapse to attending trials with the expectation that there will be less attention lapse trials overall. Piloting and calibration will be required to ensure suitable implementation of this control manipulation. Importantly, another potential approach for evaluating intervention efficacy could use a between-subjects design, such as a control group receiving sham feedback yoked to matched participants (i.e., along key demographic characteristics such as age, gender, and handedness) from a separate experimental group. This approach would require data collection from more participants than the foregoing within-subjects approach, in addition to not allowing for direct comparisons of individual differences for the intervention effect. Planned comparisons aim to test whether there are no or lesser performance effects in the control condition (i.e., non-attention lapse trials) compared to the experimental condition, even if controls show increased attentional sensitivity in response to the sham attention orienting cues (see \cite{megan2015closed}). Additionally, subsequent memory should be superior on lapse trials followed by attention reorienting compared to lapse trials without attention reorienting.

Although this approach offers predictive power that would account for individual differences in pupil and behavioral signatures, there exist limitations that need to be addressed. First, the probability that an individual gets a particular trial correct or incorrect is not solely determined by attentional state. It is possible that the individual failed to adequately encode the to-be-learned information in the first place, and thus would have failed to retrieve that information regardless of whether sustained attention was optimal or suboptimal. Second, this method requires computational resources that could efficiently compute dynamically updating model fits and subsequently control the task attention reorienting cues in real-time without sacrificing the temporal fidelity of the paradigm. Given the general accessibility of high performance computing necessary to accomplish this feat, this is less of a concern today as it would have been a decade or two ago. Nonetheless, this approach, as with the others proposed, would need to be validated in the context of computational performance \emph{and} task-based memory performance and accuracy.

\subsection{Concept 3: Momentary Millisecond Threshold-Based Pupil Diameter Method}
\label{section:concept-3}
The following two methods move away from slow rhythm attention signals towards a ms-by-ms fast rhythm attention signature by leveraging measures of significant deviations in the continuous time series of pupil diameter in the seconds directly before each retrieval probe; hence, no signal from previous trials is taken into account in this approach (see Figure \ref{fig:key-figure-ms-by-ms}a). Specifically, \emph{Concept 3} described here will be similar to \emph{Concept 1} (Section \ref{section:concept-1}) in that there will be no training phase as described in the predictive modeling approach with machine learning based tools. Similar to before, two potential approaches can be utilized to determine the threshold for which to trigger attention orienting cues during moments of attention lapsing as determined by pupil diameter signatures that deviate outside the bounds of this threshold. Determining which thresholding approach is most appropriate for accurately cuing participants' attention lapses will need to be determined by trial-and-error piloting with iterative testing prior to full scale implementation in an experimental intervention setup.

Unlike the other approaches described previously, this ms-by-ms approach raises a new concern as a function of the pupillometry methods being implemented. Stimulus onset and changes in stimulus presentation on the computer screen during a task using pupillometry require measurement of baseline pupil diameter to account for changes in luminance (as changing the pixels on the screen will result in rapid changes in luminance --- and subsequently pupil diameter --- which results in noise when using pupil as a signal for lapses in sustained attention).\footnote{One alternative design to avoid potential luminance effects is the use of auditory memory tasks to intervene on. In this scenario, participants would continue to perform the auditory memory task by providing responses via the keyboard. To control for luminance changes, participants would --- for instance --- continuously view a static text box (in the center of the screen) that would match the luminance of an attention orienting cue, which would appear when in-the-moment pupil deviations exceed the specified threshold. Though, one potential concern with an auditory stimulus approach are spikes in pupil signal as an artifact of being startled by noise-stimulus onset; this would need to be carefully tested and calibrated to assess possible negative consequences of using auditory cues over visual cues. Overall, differences in intervention effects between auditory and visual memory tasks should be assessed in future studies once the real-time closed-loop pupillometry methodology is designed and validated.} Using a 2-s period to assess baseline pupil diameter before each trial will help minimize noise in pupil diameter signal spikes from sudden changes in luminance due to novel stimulus onset (\cite{unsworth2017importance}). Specifically, prior works have shown that to accurately assay trial-level tonic lapses in attention, pupil diameter should be computed from epochs from 1-s pre-retrieval to 0-s probe onset (\cite{madore2020memory, unsworth2016pupillary, unsworth2017importance, unsworth2018tracking}). This real-time, ms-by-ms triggering method described here will thus need to account for rapid changes in pupil diameter as a function of stimulus onset. Data will need to be pre-processed, in real-time, to filter out this artifact fast enough so the intervention is not lagged by the in-the-moment computational processing power needed to both (a) pre-process the noisy pupil data, and (b) make a decision about whether or not to present an attention reorienting cue to the participant.

Importantly, RT is not considered in this approach as there would be no RT signal available for the algorithm to use; to clarify, this is because the real-time attention reorienting trigger is based on the 1-s directly before the actual retrieval task (where behavior --- here being RT --- has not yet occurred). Though, in Sections \ref{section:concept-1} and \ref{section:concept-2}, RT can technically be assayed because these methodological concepts rely on the data from the entire trial(s) directly preceding the current trial.

\begin{figure}[ht!] 
    \centering
    \includegraphics[width=1\linewidth]{figs/new_model_trial-by-trial.png}
    \includegraphics[width=1\linewidth]{figs/model_ms-by-ms-label.png}
    \caption{\textbf{\textit{a,}} Threshold-based in-the-moment immediate attention lapses with pupillometry. \textbf{\textit{b,}} Machine learning predictive in-the-moment immediate attention lapses with pupillometry. Created with \href{BioRender.com}{BioRender.com}.}
    \label{fig:key-figure-ms-by-ms}
\end{figure}

\subsection{Concept 4: Momentary Millisecond Predictive Model-Based Pupil Diameter Method}
\label{section:concept-4}
This last method utilizes the predictive modeling approach trained on data from the relevant epochs before each retrieval probe described in Section \ref{section:concept-3} (see Figure \ref{fig:key-figure-ms-by-ms}b). Similarly, this predictive model would only use tonic pupil diameter in the 1-s epoch directly before the retrieval cue, and thus RT would not yet be available during the actual intervention task since the specific trial would not yet have occurred until after the fast rhythm measure of attention is assayed. Though, like Section \ref{section:concept-2}, the model would predict the binary outcome of correct (1) or incorrect (0) as a function of pupil diameter, except in this instance the pupil diameter in the 1-s epoch directly before the retrieval cue would be the predictor for the outcome of memory performance on that same trial (in Section \ref{section:concept-2} the predictive model used the signal from the \emph{n - 1} trial with respect to the current trial). This predictive modeling approach once again offers a more fine-tuned, individualized approach to predicting memory failures as a result of in-the-moment lapses of attention; however, it is still limited by the potential influence of training effects that could confound performance during the intervention task and/or the mechanism underlying the measured behavior. 

Uniquely, given that this approach only relies on pupil signatures from 1-s epochs directly before the retrieval probe, a lot more data will need to be collected over more trials than the predictive approach described in Section \ref{section:concept-2}. This is necessary in order to have similar amounts of training measures compared to what would be collected over the course of a trial-to-trial based-method, such as the \emph{n - 3} trial predictive approach described previously in Section \ref{section:concept-1}. Practically, acquiring these training data may be more time consuming per participant and may also have even more pronounced consequential influences on the subsequent intervention phase given that participants would spend much more time in the training phase than in that of the predictive modeling method framework described in Section \ref{section:concept-2}. Ideally, one should aim for this training phase data collection period to be the same amount of time, or less time, than the subsequent intervention version of the task --- not more.

\section{Limitations and Future Directions}
\subsection{Reactive Training Effects}
In both predictive model-based concepts, there is a potential concern for adaptation and/or reactive effects from the first round of data collection for model training, which may have downstream impacts on the mechanisms guiding memory performance on the intervention-version of the trial. Careful testing of these possible performance effects from Time 1 (training) to Time 2 (intervention) is necessary to weigh the potential pros and cons of the respective closed-loop method concepts described previously.

\subsection{Pupillary Data Standardization}
In all of these method concepts, pupil measures should be \emph{z}-scored for post hoc between-subjects analyses of pupil diameter to account for individual variability in the units of measure for tonic pupil diameter. The pupil variability signal is an arbitrary unit with respect to the distance the individual is from the computer screen and the distance in which their eyes are recessed behind the eye orbit (socket). A functional setup following best practices should hold the distance from the chin rest to the eye tracker/computer screen constant across all experimental runs. Though, there will still be a small noise signal due to individual variability in the position of eyes within the socket. Applying a \emph{z}-score data transformation can supply more confidence when conducting post hoc between-subjects analyses of pupil variability across individuals.

\subsection{Control Condition with Sham Feedback}
Each of the four conceptual methods proposed above would ultimately need to be validated by 2 × 2 experimental designs. Specifically, the first independent variable in such a 2 × 2 design matrix would be attention state (\emph{attending, lapsing}), while the second would be feedback type (\emph{participant-driven, sham}). The control condition could either receive within-subject sham feedback that occurs on a predefined random subset of trials (i.e., not driven at all by attention state on those trials), or even a between-subjects sham feedback yoked to participants matched along key demographic characteristics --- from data collected in a between-subjects intervention-only condition --- such that feedback would be based on other individuals’ pupils (similar to \cite{megan2015closed}). The efficacy of the attention reorienting cue intervention within this validation setup would expectedly reveal significant boosts in memory performance for individuals in the experimental feedback condition relative to the control group receiving sham feedback.

\subsection{Ethical Considerations}
Lastly, researchers must carefully consider potential negative impacts on participants with neurofeedback implementations. To illustrate, participants should have enough information about the study to provide their informed consent about risks and benefits, or withdraw from participation in the study (\cite{cherkaoui2021ethical}). If participants are not made aware that their cognition will be causally manipulated as a function of their own brain and behavioral signals, this fundamentally opposes the ethical notion of cognitive liberty --- or an individual's right to have full autonomy over their mind and brain chemistry (\cite{lavazza2018freedom, bublitz2013my, sententia2004neuroethical}). Thus, despite the transformative potential that NF interventions can have on illustrating causal links between brain and behavior (\cite{cherkaoui2021ethical}), researchers must weigh the ethical costs of withholding critical information about a NF intervention with potential limitations on the interpretation of results if participants were explicitly made aware of the NF signal being used to manipulate behavior.

\section{Conclusion}
Here, I presented a novel conceptual framework for feasible closed-loop interventions of attention lapsing on subsequent memory using tonic pupil diameter biofeedback. Specifically, I outlined four potential methodological configurations --- each with their own strengths and limitations --- that could be designed and implemented for studies of attention reorientation on forms of memory with the specific aim of illustrating causal links between attention lapsing and memory performance. These implementations unlock numerous, novel fruitful avenues of research on a variety of populations that may not be traditionally well suited for fMRI closed-loop approaches. Importantly, if real-time manipulations of attention using each participant's own pupil signal for reorienting causes substantial effects on memory performance, then this closed-loop intervention will reveal transformative potential for both basic science and translational aims by (i) elucidating causal links between attention and memory, and (ii) showcasing a feasible, non-invasive attention-training implementation potentially useful for a wide gamut of individuals spanning diversity across lifespan, psychopathology, location, and even ecologically valid scenarios. Though, closed-loop biofeedback with pupillometry does not only need to be used as a training regime for individuals to learn to reorient their attention during crucial preparatory periods for learning information in-the-moment, but can also extend to longer forms of memory during retrieval events where the fidelity of memory reinstatement is threatened by environmental cues adversarial to the episodic retrieval task at hand. In sum, the possibilities of these approaches for basic science and translational aims span far beyond those described here, and will ultimately move human memory research into a new frontier.